
\subsection{滤过 \texorpdfstring{$A_\infty$}{A-infinity} 双模
\texorpdfstring{$C(L^{(1)},L^{(0)};\Lambda_{0,\mathrm{nov}})$}
{C(L1,L0;Lambda0nov)}}
\label{subsec:3-7-4}

本小节的主要结果是定理 3.7.21。本小节取 $R=\mathbb Q$。

设 $L^{(0)},L^{(1)}$ 是 $M$ 中的一对相对 spin 拉格朗日子流形。我们先定义
一个 $\Lambda_{0,\mathrm{nov}}$-模
$C(L^{(1)},L^{(0)};\Lambda_{0,\mathrm{nov}})$，并在其上构造滤过
$A_\infty$ 双模结构。为此并不需要假设 $L^{(i)}$ 无障碍。本小节假设
$L^{(0)}$ 与 $L^{(1)}$ 彼此横截；下一小节将处理二者干净相交的情形。

令
\[
  \bigl(C(L^{(0)};\Lambda_{0,\mathrm{nov}}),\mathfrak m^{(0)}\bigr),
  \qquad
  \bigl(C(L^{(1)};\Lambda_{0,\mathrm{nov}}),\mathfrak m^{(1)}\bigr)
\]
为定理 3.5.11 所构造的滤过 $A_\infty$ 代数。定义
\begin{align*}
 CF(L^{(1)},L^{(0)};\ell_0)
 &=CF(L^{(1)},L^{(0)};\ell_0;\mathbb Q)\\
 &:=\widehat{\bigoplus}_{
   [\ell_p,w]\in\operatorname{Cr}(L^{(0)},L^{(1)};\ell_0)}
   \mathbb Q[\ell_p,w],\\
 CF(L^{(1)},L^{(0)})
 &:=\widehat{\bigoplus}_{
   [\ell_0]\in\pi_0(\Omega(L^{(0)},L^{(1)}))}
   CF(L^{(1)},L^{(0)};\ell_0)\\
 &=\widehat{\bigoplus}_{
   [\ell_p,w]\in\operatorname{Cr}(L^{(0)},L^{(1)})}
   \mathbb Q[\ell_p,w].
\end{align*}
这里 $\operatorname{Cr}(L^{(0)},L^{(1)};\ell_0)$ 与
$\operatorname{Cr}(L^{(0)},L^{(1)})$ 是作用量泛函 $\mathcal A$ 的临界点
等价类 $[\ell_p,w]$ 之集（见第 2.2.2 小节）；上式第二个直和遍历
$\Omega(L^{(0)},L^{(1)})$ 的全部连通分支
\translatornote{原书第二行把分支指标排成
$\ell_0\in\pi_0(\Omega(L^{(1)},\Omega(L^{(0)})))$，其中道路空间的第二个
自变量被误排成另一个道路空间，且次序与紧接着的说明和全章定义相反。译文依
$\Omega(L^{(0)},L^{(1)})$ 的连通分支改正，并以 $[\ell_0]$ 明示分支类。}。符号
$\widehat{\bigoplus}$ 表示直和的完备化。$\Gamma$-等价关系见定义 2.2.4，
而 $\ell_p$ 是通向交点 $p\in L^{(1)}\cap L^{(0)}$ 的常值道路。

定义
\[
 \widehat{CF}(L^{(1)},L^{(0)};\Lambda_{\mathrm{nov}})
 :=CF(L^{(1)},L^{(0)})
   \widehat\otimes_{\mathbb Q}\Lambda_{\mathrm{nov}},
\]
并在其上定义等价关系 $\sim$ 如下。对其中两个元素，规定
\[
 T^\lambda e^\mu[\ell_p,w]\sim
 T^{\lambda'}e^{\mu'}[\ell_{p'},w']
\]
当且仅当下列条件成立：
\begin{align}
 p&=p',
 \tag{3.7.19.1}\label{eq:3-7-19-1}\\
 \lambda+\int_w\omega
 &=\lambda'+\int_{w'}\omega,
 \tag{3.7.19.2}\label{eq:3-7-19-2}\\
 2\mu+\mu([\ell_p,w])
 &=2\mu'+\mu([\ell_{p'},w']).
 \tag{3.7.19.3}\label{eq:3-7-19-3}
\end{align}
这里 $\mu([\ell_p,w])$ 是第 2.2.2 小节给出的马斯洛夫--莫尔斯指标。
这些条件与定义 2.2.4 的 $\Gamma$-等价
\translatornote{原书此处把“$\Gamma$-等价”的出处写成定义 2.2.12；
该定义实际定义马斯洛夫--莫尔斯指标，而 $\Gamma$-等价定义于 2.2.4。
译文分别恢复两个对象的正确出处。}相容，因而 $\sim$ 确实在
$\widehat{CF}(L^{(1)},L^{(0)};\Lambda_{\mathrm{nov}})$ 上定义了等价关系。

再在
\[
 \widehat{CF}(L^{(1)},L^{(0)};\Lambda_{\mathrm{nov}})
\]
上定义能量滤过。对其中的元素 $T^\lambda e^\mu[\ell_p,w]$，规定
\[
 T^\lambda e^\mu[\ell_p,w]\in
 F^{\lambda'}
 \widehat{CF}(L^{(1)},L^{(0)};\Lambda_{0,\mathrm{nov}})
 \quad\Longleftrightarrow\quad
 \lambda+\int_w\omega\geq\lambda'.
\]
这个滤过显然在商空间
$\widehat{CF}(L^{(1)},L^{(0)};\Lambda_{\mathrm{nov}})/\sim$ 上诱导滤过。

\begin{definition}{3.7.20}\label{def:3-7-20}
以 $C(L^{(1)},L^{(0)};\Lambda_{0,\mathrm{nov}})$ 表示
$\widehat{CF}(L^{(1)},L^{(0)};\Lambda_{\mathrm{nov}})/\sim$
关于能量滤过作完备化后所得的非负能量部分。也就是说，
\[
 C(L^{(1)},L^{(0)};\Lambda_{0,\mathrm{nov}})
 =F^0CF(L^{(1)},L^{(0)};\Lambda_{\mathrm{nov}}).
\]
\end{definition}

把元素 $T^\lambda e^\mu[\ell_p,w]$ 的次数定义为
$2\mu+\mu([\ell_p,w])$，便得到一个滤过分次自由
$\Lambda_{0,\mathrm{nov}}$-模。容易看出，
$C(L^{(1)},L^{(0)};\Lambda_{0,\mathrm{nov}})$ 满足
(3.2.12.1)--(3.2.12.5)。还要注意，作为
$\Lambda_{0,\mathrm{nov}}$-模，它同以交点生成的自由模之滤过完备化同构：
\[
 C(L^{(1)},L^{(0)};\Lambda_{0,\mathrm{nov}})
 \cong
 \widehat{\bigoplus}_{p\in L^{(1)}\cap L^{(0)}}
 \Lambda_{0,\mathrm{nov}}[p].
\]

仍令 $\{J_t\}_t=\{J_t\}_{0\leq t\leq1}$ 为前面所用的依赖于 $t$ 的殆复结构族。
下面是本小节的主要定理。

\begin{theorem}{3.7.21}\label{thm:3-7-21}
设 $(L^{(0)},L^{(1)})$ 是一对彼此横截的相对 spin 拉格朗日子流形。则
$C(L^{(1)},L^{(0)};\Lambda_{0,\mathrm{nov}})$ 上存在
\[
 \begin{gathered}
 \text{左 }\bigl(C(L^{(1)};\Lambda_{0,\mathrm{nov}}),\mathfrak m^{(1)}\bigr),\\
 \text{右 }\bigl(C(L^{(0)};\Lambda_{0,\mathrm{nov}}),\mathfrak m^{(0)}\bigr)
 \end{gathered}
\]
的滤过 $A_\infty$ 双模结构\liftIndex{I116}，并且它是
$G(L^{(1)},L^{(0)})$-有隙的。此外，上述两个代数的同伦幺元之对
$\{\mathbf e_1,\mathbf e_0\}$ 作为同伦幺元起作用。
\end{theorem}

$G(L^{(1)},L^{(0)})$ 将在证明中由式 \eqref{eq:3-7-43} 定义。

\begin{proof}
横截性问题需要像定理 3.5.11 那样精确处理；这一部分的证明只是第 3.5 节论证的
小改动。同伦幺元的讨论与第 7.3 节类似，符号见第 8.10 节。

取
\[
 [\ell_p,w_1],[\ell_q,w_2]
 \in\operatorname{Cr}(L^{(0)},L^{(1)}).
\]
为在 $C(L^{(1)},L^{(0)};\Lambda_{0,\mathrm{nov}})$ 上构造滤过
$A_\infty$ 双模结构，需要定义带标记稳定断裂弗洛尔轨迹的模空间
\liftIndex{I127}
\[
 \mathcal M_{k_1,k_0}([\ell_p,w_1],[\ell_q,w_2])
 :=
 \mathcal M_{k_1,k_0}
 (L^{(1)},L^{(0)};[\ell_p,w_1],[\ell_q,w_2]).
\]
这里各参数的次序已经固定。特例
$\mathcal M_{0,0}([\ell_p,w_1],[\ell_q,w_2])$
与第 2.4.2 小节定义的
$\mathcal M([\ell_p,w_1],[\ell_q,w_2];\{J_t\}_t)$ 一致。

令 $\Sigma$ 为 $\mathbb R\times[0,1]$，并允许在
$\mathbb R\times(0,1)$ 中有限多个点处附着有限层球面分支。球面气泡的纤维积处理
与 \textbf{[FuOn99II]} 相同，故不再重复。令 $u\colon\Sigma\to M$ 光滑，并取
\[
 \tau_j^{(0)}\in\mathbb R\ (j=1,\ldots,k_0),\qquad
 \tau_j^{(1)}\in\mathbb R\ (j=1,\ldots,k_1).
\]
把这些参数记为
\[
 \vec\tau^{\,(i)}
 =(\tau_1^{(i)},\ldots,\tau_{k_i}^{(i)})\quad(i=0,1),
 \qquad
 \vec\tau=(\vec\tau^{\,(0)},\vec\tau^{\,(1)}).
\]
我们要求它们满足下列条件（见图 3.7.1）。

\begin{figure}[htbp]
  \centering
  \setlength{\unitlength}{0.9mm}
  \begin{picture}(126,70)
    \qbezier(8,34)(63,69)(118,34)
    \qbezier(8,34)(63,-1)(118,34)
    \put(5,31){\line(1,1){7}}
    \put(5,37){\line(1,-1){7}}
    \put(114,31){\line(1,1){7}}
    \put(114,37){\line(1,-1){7}}
    \put(27,46){\circle*{1.6}}
    \put(44,55){\circle*{1.6}}
    \put(82,55){\circle*{1.6}}
    \put(101,46){\circle*{1.6}}
    \put(27,18){\circle*{1.6}}
    \put(98,20){\circle*{1.6}}
    \put(15,51){$\tau_{k_1}^{(1)}$}
    \put(98,51){$\tau_1^{(1)}$}
    \put(17,8){$\tau_1^{(0)}$}
    \put(92,9){$\tau_{k_0}^{(0)}$}
    \put(119,43){$L^{(0)}$}
    \put(119,24){$L^{(1)}$}
    \put(48,53){\circle{17}}
    \put(48,66){\circle{12}}
    \put(77,53){\circle{17}}
    \put(77,66){\circle{12}}
    \qbezier[12](53,68)(63,74)(72,68)
    \qbezier[12](38,9)(63,1)(88,9)
  \end{picture}
  \par\smallskip
  \renewcommand{\thefigure}{3.7.1}
  \refstepcounter{figure}
  \textbf{\figurename~\thefigure}\label{fig:3-7-1}
\end{figure}

\begin{equation}
\begin{gathered}
 \frac{\partial u}{\partial\tau}
 +J_t\frac{\partial u}{\partial t}=0
 \quad\text{在 }\mathbb R\times[0,1]\text{ 上};\\
 \text{若一棵球面气泡树扎根于 }(\tau,t),
 \text{ 则树中各球面均为 }J_t\text{-全纯的}.
\end{gathered}
\tag{3.7.22.1}\label{eq:3-7-22-1}
\end{equation}
\begin{equation}
 u(\mathbb R\times\{0\})\subset L^{(0)},\qquad
 u(\mathbb R\times\{1\})\subset L^{(1)}.
\tag{3.7.22.2}\label{eq:3-7-22-2}
\end{equation}
\begin{equation}
 \lim_{\tau\to-\infty}u(\tau,t)=p,\qquad
 \lim_{\tau\to+\infty}u(\tau,t)=q.
\tag{3.7.22.3}\label{eq:3-7-22-3}
\end{equation}
\begin{equation}
\begin{gathered}
 \tau_j^{(0)}<\tau_{j+1}^{(0)}
 \quad(j=1,\ldots,k_0-1),\\
 \tau_j^{(1)}>\tau_{j+1}^{(1)}
 \quad(j=1,\ldots,k_1-1).
\end{gathered}
\tag{3.7.22.4}\label{eq:3-7-22-4}
\end{equation}
\begin{equation}
 w_1\#u\sim w_2.
\tag{3.7.22.5}\label{eq:3-7-22-5}
\end{equation}
\begin{equation}
\begin{gathered}
 ((\Sigma,\vec\tau),u)\text{ 是稳定的，即满足}\\
 \bigl\{\varphi\colon\Sigma\to\Sigma\mid
 \varphi(\tau_j^{(i)})=\tau_j^{(i)},\
 u\circ\varphi=u\bigr\}\text{ 是有限集}.
\end{gathered}
\tag{3.7.22.6}\label{eq:3-7-22-6}
\end{equation}
在 \eqref{eq:3-7-22-5} 中，等价关系 $\sim$ 采用定义 2.2.4 的记号；
$w_1\#u$ 是 $w_1$ 与 $u$ 的自然拼接（见第 2.2.1 小节）。

\begin{remark}{3.7.23}\label{rem:3-7-23}
\begin{enumerate}[label=\textup{(\arabic*)},leftmargin=2.8em]
 \item 我们赋予
 $C(L^{(1)},L^{(0)};\Lambda_{0,\mathrm{nov}})$ 如下滤过
 $A_\infty$ 双模结构：左作用来自
 \[
  \bigl(C(L^{(1)};\Lambda_{0,\mathrm{nov}}),\mathfrak m^{(1)}\bigr),
 \]
 右作用来自
 \[
  \bigl(C(L^{(0)};\Lambda_{0,\mathrm{nov}}),\mathfrak m^{(0)}\bigr).
 \]
 乍看之下，似乎反过来取右 $L^{(1)}$、左 $L^{(0)}$ 更自然。
 定理 3.7.21 采用上述左右次序，是因为我们考虑从
 $L^{(0)}$ 到 $L^{(1)}$ 的道路空间，所以
 $u\colon\Sigma\to M$ 的自然边界条件正是 \eqref{eq:3-7-22-2}。
 见图 3.7.2。

\begin{figure}[htbp]
  \centering
  \setlength{\unitlength}{0.82mm}
  \begin{picture}(132,91)
    \put(10,65){\line(1,0){112}}
    \put(10,48){\line(1,0){112}}
    \multiput(16,48)(15,0){7}{\line(1,2){8.5}}
    \put(5,55){$z_1\leftarrow$}
    \put(119,55){$\rightarrow z_0$}
    \put(28,63){\line(1,1){3}}
    \put(28,66){\line(1,-1){3}}
    \put(54,63){\line(1,1){3}}
    \put(54,66){\line(1,-1){3}}
    \put(91,63){\line(1,1){3}}
    \put(91,66){\line(1,-1){3}}
    \put(106,63){\line(1,1){3}}
    \put(106,66){\line(1,-1){3}}
    \put(24,71){$\tau_{k_1}^{(1)}=z_{k_1+1}$}
    \put(100,71){$\tau_1^{(1)}=z_2$}
    \put(23,46){\line(1,1){3}}
    \put(23,49){\line(1,-1){3}}
    \put(47,46){\line(1,1){3}}
    \put(47,49){\line(1,-1){3}}
    \put(83,46){\line(1,1){3}}
    \put(83,49){\line(1,-1){3}}
    \put(111,46){\line(1,1){3}}
    \put(111,49){\line(1,-1){3}}
    \put(16,38){$\tau_1^{(0)}=z_{k_1+2}$}
    \put(93,38){$\tau_{k_0}^{(0)}=z_{k_0+k_1+1}$}
    \put(66,17){\circle{34}}
    \put(49,17){\circle*{1.5}}
    \put(83,17){\circle*{1.5}}
    \put(52,28){\line(1,1){2.5}}
    \put(52,30.5){\line(1,-1){2.5}}
    \put(77.5,28){\line(1,1){2.5}}
    \put(77.5,30.5){\line(1,-1){2.5}}
    \put(55,2){\line(1,1){2.5}}
    \put(55,4.5){\line(1,-1){2.5}}
    \put(76,3){\line(1,1){2.5}}
    \put(76,5.5){\line(1,-1){2.5}}
    \put(43,16){$z_1$}
    \put(85,16){$z_0$}
    \put(45,31){$z_{k_1+1}$}
    \put(78,31){$z_2$}
    \put(45,-1){$z_{k_1+2}$}
    \put(78,-1){$z_{k_0+k_1+1}$}
  \end{picture}
  \par\smallskip
  \renewcommand{\thefigure}{3.7.2}
  \refstepcounter{figure}
  \textbf{\figurename~\thefigure}\label{fig:3-7-2}
\end{figure}

 圆盘作为有边曲面时，其边界的标准定向是逆时针方向。因此，标记点
 $(\tau_j^{(i)},i)$（$i=0,1$，$j=1,\ldots,k_i$）与 $-\infty$ 的自然次序为
 \[
  (\tau_1^{(1)},1),\ldots,(\tau_{k_1}^{(1)},1),-\infty,
  (\tau_1^{(0)},0),\ldots,(\tau_{k_0}^{(0)},0).
 \]
 这说明 $L^{(1)}$ 上的输入排在前面，其后是 $[\ell_p,w_1]$，最后才是
 $L^{(0)}$ 上的输入。

 \item 在把 $\mathbb R\times[0,1]$ 共形等价为圆盘时，$-\infty$ 与
 $+\infty$ 分别对应 $-1$ 与 $+1$。在后面的论证中，这两点将分别对应
 第 1 与第 0 个标记点 $z_1,z_0$。见第 8.7、8.8 节。边界上标记点的
 次序在讨论定向问题时十分重要。我们希望继续以 $z_0,z_1$
 分别表示 $+1,-1$。第 8.1、8.8 节采用
 \[
 \begin{gathered}
  z_2=(\tau_1^{(1)},1),\ldots,
  z_{k_1+1}=(\tau_{k_1}^{(1)},1),\\
  z_{k_1+2}=(\tau_1^{(0)},0),\ldots,
  z_{k_0+k_1+1}=(\tau_{k_0}^{(0)},0);
 \end{gathered}
 \]
 另见图 3.7.2。
 \translatornote{原书第二行首项写成 $z_{k_0+2}$；它既不承接前一项
 $z_{k_1+1}$，也与图 3.7.2 的标号不符。译文依连续编号和原图改为
 $z_{k_1+2}$。}
\end{enumerate}
\end{remark}

考虑 $\mathbb R$-作用
\[
 u(\mathord\cdot,\mathord\cdot)
 \longmapsto u(\mathord\cdot+\tau_0,\mathord\cdot),
 \qquad \tau_0\in\mathbb R.
\]
它在式 \eqref{eq:3-7-22-1}--\eqref{eq:3-7-22-6} 的解集上诱导自然等价关系
$\sim_{\mathbb R}$。

\begin{definition}{3.7.24}\label{def:3-7-24}
以
\[
 \mathring{\mathcal M}_{k_1,k_0}([\ell_p,w_1],[\ell_q,w_2])
 =
 \mathring{\mathcal M}_{k_1,k_0}
 (L^{(1)},L^{(0)};[\ell_p,w_1],[\ell_q,w_2])
\]
表示满足 \eqref{eq:3-7-22-1}--\eqref{eq:3-7-22-6} 的
$((\Sigma,\vec\tau),u)$ 的全部 $\sim_{\mathbb R}$-等价类。
\end{definition}

评价映射为
\begin{equation}
\begin{aligned}
 \operatorname{ev}
 &=(\operatorname{ev}^{(1)},\operatorname{ev}^{(0)})\\
 &=\bigl((\operatorname{ev}^{(1)}_1,\ldots,
   \operatorname{ev}^{(1)}_{k_1}),
   (\operatorname{ev}^{(0)}_1,\ldots,
   \operatorname{ev}^{(0)}_{k_0})\bigr)\\
 &\colon
 \mathring{\mathcal M}_{k_1,k_0}([\ell_p,w],[\ell_q,w'])
 \longrightarrow (L^{(1)})^{k_1}\times(L^{(0)})^{k_0},
\end{aligned}
\tag{3.7.25}\label{eq:3-7-25}
\end{equation}
其中
\[
 \operatorname{ev}^{(i)}_j((\Sigma,\vec\tau),u)
 =u(\tau_j^{(i)},i)\in L^{(i)},\qquad i=0,1.
\]

为陈述下一个命题，先说明记号。回忆命题 3.4.2 已对
$\beta\in\Pi(M;L)$ 在 $\mathcal M_k(\beta)$ 上构造仓西结构
（$\Pi(M;L)$ 定义于 2.4.17）。模空间
$\mathcal M_k^{\mathrm{main}}(\beta)$ 及其仓西结构依赖于 $L,\beta$
以及相容殆复结构 $J$ 的选择。为标明这些选择，写作
\[
 \mathcal M_k^{\mathrm{main}}(L;\beta;J).
\]

\liftIndexRangeStart{I126}
\begin{proposition}{3.7.26}\label{prop:3-7-26}
设 $(L^{(0)},L^{(1)})$ 是一对彼此横截的相对 spin 拉格朗日子流形，并且
\[
 [\ell_p,w_1],[\ell_q,w_2]
 \in\operatorname{Cr}(L^{(0)},L^{(1)}).
\]
则
\[
 \mathring{\mathcal M}_{k_1,k_0}
 (L^{(1)},L^{(0)};[\ell_p,w_1],[\ell_q,w_2])
\]
有一个紧化，记作
\[
 \mathcal M_{k_1,k_0}
 (L^{(1)},L^{(0)};[\ell_p,w_1],[\ell_q,w_2]).
\]
评价映射 \eqref{eq:3-7-25} 延拓到这个紧化，并且是弱浸没、强连续且光滑的。

该紧化带有定向仓西结构；这些结构彼此相容，并与命题 3.4.2 所构造的
\[
 \mathcal M_k^{\mathrm{main}}(L^{(0)};\beta^{(0)};J_0),
 \qquad
\mathcal M_k^{\mathrm{main}}(L^{(1)};\beta^{(1)};J_1)
\]
上的仓西结构相容
\translatornote{原书命题把两端生成元写成
$\operatorname{Cr}(L^{(1)},L^{(0)})$ 的元素；本小节此前所有生成元和道路空间
均按 $\operatorname{Cr}(L^{(0)},L^{(1)})$ 定义，命题所紧化的也是同一个模空间。
译文据此恢复既定次序。}。
\end{proposition}

命题 3.7.26 中的相容性，是指在模空间端部通过纤维积所表达的相容性
（见第 7.1.1 小节）。

下面只描述紧化
$\mathcal M_{k_1,k_0}
(L^{(1)},L^{(0)};[\ell_p,w_1],[\ell_q,w_2])$；
其仓西结构的构造与第 7.1 节讨论的其他情形相同。为描述这个紧化，先准备一些记号。

\begin{definition}{3.7.27}\label{def:3-7-27}
令
\[
 l_0,l_1,m_0,m_1,\quad
 m_{0,1},\ldots,m_{0,l_0},\quad
 m_{1,1},\ldots,m_{1,l_1}
\]
为非负整数，并满足
\[
 k_i=m_i+\sum_{a=1}^{l_i}m_{i,a},\qquad i=0,1.
\]
取子集
\[
 A_i\subset\{1,\ldots,m_i+l_i\},\qquad \#A_i=l_i,
\]
并写成
\[
 A_i=\{\sigma_i(1),\ldots,\sigma_i(l_i)\},
 \qquad \sigma_i(a)<\sigma_i(a+1).
\]
令
\[
 \pi_{i,1}\colon(L^{(i)})^{m_i+l_i}\longrightarrow(L^{(i)})^{l_i}
\]
为投影
\[
 \pi_{i,1}(x_1,\ldots,x_{m_i+l_i})
 =(x_{\sigma_i(1)},\ldots,x_{\sigma_i(l_i)}),
\]
并令
\[
\pi_{i,2}\colon(L^{(i)})^{m_i+l_i}\longrightarrow(L^{(i)})^{m_i}
\]
为投到其余因子的投影
\translatornote{原书把两个投影的源写成
$(L^{(i)})^{k_i+l_i}$，并把第二个投影的目标写成
$(L^{(i)})^{k_i}$。然而 $A_i$ 是
$\{1,\ldots,m_i+l_i\}$ 的子集，且紧接的 (3.7.28) 要把这些投影同主分支
$\mathring{\mathcal M}_{l_1+m_1,l_0+m_0}$ 的评价映射复合；因此源必须有
$m_i+l_i$ 个因子，删去 $l_i$ 个附着位置后目标有 $m_i$ 个因子。译文依此改正。}。
以 $A$ 表示上述全部数据
\[
 l_0,l_1,m_0,m_1,\{m_{0,a}\},\{m_{1,a}\},A_0,A_1.
\]
\end{definition}

据此，置
\begin{equation}
\begin{aligned}
&\mathcal M^A_{k_1,k_0}([\ell_p,w_1],[\ell_q,w_2])\\
&:=
\coprod_{\substack{
 w'_2\#\sum_{a=1}^{l_1}\beta_{1,a}
       \#\sum_{b=1}^{l_0}\beta_{0,b}=w_2}}
 \mathring{\mathcal M}_{l_1+m_1,l_0+m_0}
 ([\ell_p,w_1],[\ell_q,w'_2])\\
&\qquad{}_{(\pi_{1,1},\pi_{0,1})\circ\operatorname{ev}}
 \!\times_{\overrightarrow{\operatorname{ev}}_0}
 \left(
 \begin{gathered}
  \prod_{a=1}^{l_1}
  \mathcal M_{m_{1,a}+1}^{\mathrm{main}}
  (L^{(1)};\beta_{1,a};J_1)\\[-0.2em]
  {}\times
  \prod_{b=1}^{l_0}
  \mathcal M_{m_{0,b}+1}^{\mathrm{main}}
  (L^{(0)};\beta_{0,b};J_0)
 \end{gathered}
 \right),
\end{aligned}
\tag{3.7.28}\label{eq:3-7-28}
\end{equation}
其中
$\overrightarrow{\operatorname{ev}}_0
=(\operatorname{ev}_0,\ldots,\operatorname{ev}_0)$，并定义
\begin{equation}
 \mathcal M^{\mathrm{unbr}}_{k_1,k_0}
 ([\ell_p,w_1],[\ell_q,w_2])
 :=
 \bigcup_A
 \mathcal M^A_{k_1,k_0}
 ([\ell_p,w_1],[\ell_q,w_2]).
\tag{3.7.29}\label{eq:3-7-29}
\end{equation}
这里 \texttt{unbr} 是 \texttt{unbroken} 的缩写。
$\mathcal M^{\mathrm{unbr}}_{k_1,k_0}
([\ell_p,w_1],[\ell_q,w_2])$ 的元素称为\emph{边界带标记点的未断裂
弗洛尔轨迹}。见图 3.7.3。图中的
$\tau_j^{(0)},\tau_j^{(1)}$ 是式 \eqref{eq:3-7-28} 中元素的因子
$\mathring{\mathcal M}_{l_1+m_1,l_0+m_0}
([\ell_p,w_1],[\ell_q,w'_2])$ 上标记点的位置。

\begin{figure}[htbp]
  \centering
  \setlength{\unitlength}{0.84mm}
  \begin{picture}(132,76)
    \put(8,46){\line(1,0){116}}
    \put(8,28){\line(1,0){116}}
    \multiput(14,28)(15,0){7}{\line(1,2){9}}
    \put(14,44){\line(1,1){3}}
    \put(14,47){\line(1,-1){3}}
    \put(52,44){\line(1,1){3}}
    \put(52,47){\line(1,-1){3}}
    \put(91,44){\line(1,1){3}}
    \put(91,47){\line(1,-1){3}}
    \put(111,44){\line(1,1){3}}
    \put(111,47){\line(1,-1){3}}
    \put(16,26){\line(1,1){3}}
    \put(16,29){\line(1,-1){3}}
    \put(43,26){\line(1,1){3}}
    \put(43,29){\line(1,-1){3}}
    \put(67,26){\line(1,1){3}}
    \put(67,29){\line(1,-1){3}}
    \put(91,26){\line(1,1){3}}
    \put(91,29){\line(1,-1){3}}
    \put(36,57){\circle{20}}
    \put(36,46){\circle*{1.5}}
    \put(78,58){\circle{22}}
    \put(78,46){\circle*{1.5}}
    \put(72,70){\circle{9}}
    \put(86,70){\circle{9}}
    \put(51,16){\circle{20}}
    \put(51,28){\circle*{1.5}}
    \put(51,4){\circle{8}}
    \put(105,16){\circle{21}}
    \put(105,28){\circle*{1.5}}
    \put(96,3){\circle{8}}
    \put(116,4){\circle{9}}
    \put(73,56){$\beta_{1,1}$}
    \put(46,13){$\beta_{0,1}$}
    \put(100,13){$\beta_{0,2}$}
    \put(12,39){\scriptsize$\tau_{\sigma_1(2)}^{(1)}$}
    \put(68,39){\scriptsize$\tau_{\sigma_1(1)}^{(1)}$}
    \put(10,18){\scriptsize$\tau_1^{(0)}$}
    \put(35,31){\scriptsize$\tau_{\sigma_0(1)}^{(0)}$}
    \put(91,31){\scriptsize$\tau_{\sigma_0(2)}^{(0)}$}
  \end{picture}
  \par\smallskip
  \renewcommand{\thefigure}{3.7.3}
  \refstepcounter{figure}
  \textbf{\figurename~\thefigure}\label{fig:3-7-3}
\end{figure}

式 \eqref{eq:3-7-28} 中每个元素在作纤维积后，都有 $k_1$ 个映到
$L^{(1)}$ 的标记点和 $k_0$ 个映到 $L^{(0)}$ 的标记点。其中分别有
$m_1,m_0$ 个位于主分支，另有
$\sum_jm_{1,j},\sum_jm_{0,j}$ 个位于圆盘气泡上。按逆时针次序重新编号后，
得到评价映射
\begin{equation}
 \operatorname{ev}\colon
 \mathcal M^A_{k_1,k_0}([\ell_p,w_1],[\ell_q,w_2])
 \longrightarrow (L^{(1)})^{k_1}\times(L^{(0)})^{k_0}.
\tag{3.7.30}\label{eq:3-7-30}
\end{equation}

接着定义如下紧化
\translatornote{原书 (3.7.31) 的乘积写成 $a=1,\ldots,K-1$，而因子从
$p(a)$ 连到 $p(a+1)$；这会漏掉从固定输入 $p(0)=p$ 出发的第一段，并使
$k_{i,K}$ 完全未被使用。原式还没有把固定填充 $w_1,w_2$ 绑定到端点。
译文依图 3.7.4、(3.7.32.1)、(3.7.32.3) 所要求的 $K$ 段断裂链，统一改为
$a=1,\ldots,K$，第 $a$ 段从 $p(a-1)$ 连到 $p(a)$，并显式约定两端填充。}：
\begin{equation}
\begin{aligned}
&\mathcal M_{k_1,k_0}([\ell_p,w_1],[\ell_q,w_2])\\
&:=
\coprod_{\substack{
 K,\ w(\mathord\cdot),p(\mathord\cdot),\\
 k_{1,*},k_{0,*}}}
 \prod_{a=1}^{K}
 \mathcal M^{\mathrm{unbr}}_{k_{1,a},k_{0,a}}
 \left(
 \begin{gathered}
  [\ell_{p(a-1)},w(a-1)],\\[-0.2em]
  [\ell_{p(a)},w(a)]
 \end{gathered}
 \right),
\end{aligned}
\tag{3.7.31}\label{eq:3-7-31}
\end{equation}
右端不交并遍历 $K\geq1$ 以及满足下列条件的全部
$w(\mathord\cdot),p(\mathord\cdot),k_{1,*},k_{0,*}$；其中约定
\begin{align}
 p(0)&=p,\quad w(0)=w_1,\qquad
 p(K)=q,\quad w(K)=w_2,
\tag{3.7.32.1}\label{eq:3-7-32-1}\\
 [\ell_{p(a)},w(a)]
 &\in\operatorname{Cr}(L^{(0)},L^{(1)})
 \quad(1\leq a<K),
\tag{3.7.32.2}\label{eq:3-7-32-2}\\
 k_{1,1}+\cdots+k_{1,K}&=k_1,\qquad
 k_{0,1}+\cdots+k_{0,K}=k_0.
\tag{3.7.32.3}\label{eq:3-7-32-3}
\end{align}
\begin{figure}[htbp]
  \centering
  \setlength{\unitlength}{0.86mm}
  \begin{picture}(132,50)
    \qbezier(4,23)(18,43)(32,23)
    \qbezier(4,23)(18,3)(32,23)
    \qbezier(32,23)(47,42)(62,23)
    \qbezier(32,23)(47,4)(62,23)
    \qbezier(62,23)(78,42)(94,23)
    \qbezier(62,23)(78,4)(94,23)
    \qbezier(94,23)(111,45)(128,23)
    \qbezier(94,23)(111,1)(128,23)
    \put(1,20){\line(1,1){6}}
    \put(1,26){\line(1,-1){6}}
    \put(125,20){\line(1,1){6}}
    \put(125,26){\line(1,-1){6}}
    \put(18,42){\circle{11}}
    \put(13,50){\circle{6}}
    \put(24,50){\circle{6}}
    \put(23,7){\circle{9}}
    \put(78,41){\circle{9}}
    \put(111,4){\circle{9}}
    \put(31.2,22.2){\circle*{1.8}}
    \put(61.2,22.2){\circle*{1.8}}
    \put(93.2,22.2){\circle*{1.8}}
    \put(0,10){$\ell_p=\ell_{p(0)}$}
    \put(56,10){$\ell_{p(i)}$}
    \put(62,15){\vector(0,1){6.3}}
    \put(108,10){$\ell_q=\ell_{p(K)}$}
  \end{picture}
  \par\smallskip
  \renewcommand{\thefigure}{3.7.4}
  \refstepcounter{figure}
  \textbf{\figurename~\thefigure}\label{fig:3-7-4}
\end{figure}

式 \eqref{eq:3-7-30} 以显然方式诱导评价映射
\begin{equation}
 \operatorname{ev}\colon
 \mathcal M_{k_1,k_0}([\ell_p,w_1],[\ell_q,w_2])
 \longrightarrow (L^{(1)})^{k_1}\times(L^{(0)})^{k_0}.
\tag{3.7.33}\label{eq:3-7-33}
\end{equation}
与第 7.1.4 小节相同，可以在这个空间上定义拓扑，并证明它是 Hausdorff
紧空间。仓西结构的构造也与第 7.1 节的其他情形类似。模空间
\[
 \mathcal M_{k_1,k_0}
 (L^{(1)},L^{(0)};[\ell_p,w_1],[\ell_q,w_2])
\]
的定向在第 8.7 节定义。这就证明了命题 3.7.26。

现在仿照第 3.5 节，以 $S^k(L^{(0)};\mathbb Q)$（相应地，
$S^k(L^{(1)};\mathbb Q)$）表示由 $L^{(0)}$（相应地，$L^{(1)}$）上余维
$k$ 的光滑奇异单形生成的自由 $\mathbb Q$-模。考虑第 3.5 节所述的子模
\[
 C(L^{(i)};\mathbb Q)\subset S(L^{(i)};\mathbb Q),
 \qquad i=0,1,
\]
它们分别由可数集 $\mathcal X_1(L^{(i)})$ 生成。取
\[
 (P_j^{(i)},f_j^{(i)})\in\mathcal X_1(L^{(i)}),
 \qquad j=1,\ldots,k_i,
\]
并考虑纤维积
\begin{equation}
\begin{aligned}
&\mathcal M_{k_1,k_0}
 (L^{(1)},L^{(0)};[\ell_p,w_1],[\ell_q,w_2])\\[-0.2em]
&\quad{}_{\operatorname{ev}}
 \!\times_{
 f_1^{(1)}\times\cdots\times f_{k_1}^{(1)}
 \times f_1^{(0)}\times\cdots\times f_{k_0}^{(0)}}
 \left(
  \prod_{i=1}^{k_1}P_i^{(1)}
  \times\prod_{i=1}^{k_0}P_i^{(0)}
 \right).
\end{aligned}
\tag{3.7.34}\label{eq:3-7-34}
\end{equation}
记 $\vec P^{(i)}=(P_1^{(i)},\ldots,P_{k_i}^{(i)})$，并把
\eqref{eq:3-7-34} 记作
\begin{equation}
 \mathcal M_{k_1,k_0}
 (L^{(1)},L^{(0)};[\ell_p,w_1],[\ell_q,w_2];
 \vec P^{(1)},\vec P^{(0)}).
\tag{3.7.35}\label{eq:3-7-35}
\end{equation}
由命题 3.7.26，它带有仓西结构。

需要指明 $J_i$ 等选择时，我们分别写
\[
\begin{gathered}
 \mathcal M_k(L^{(i)};\beta^{(i)};J_i;\vec P^{(i)}),\\
 \mathcal M_{k_1,k_0}
 (L^{(1)},L^{(0)};[\ell_p,w_1],[\ell_q,w_2];
 \vec P^{(1)},\vec P^{(0)};\{J_t\}_t)
\end{gathered}
\]
来代替
\[
\begin{gathered}
 \mathcal M_k(\beta^{(i)};\vec P^{(i)}),\\
 \mathcal M_{k_1,k_0}
 (L^{(1)},L^{(0)};[\ell_p,w_1],[\ell_q,w_2];
 \vec P^{(1)},\vec P^{(0)}).
\end{gathered}
\]
回忆在第 3.5 节中，为了定义与 $L^{(i)}$ 相伴的滤过 $A_\infty$ 代数，
我们已经在
$\mathcal M_k(L^{(i)};\beta^{(i)};J_i;\vec P^{(i)})$
上选取多值截面系（以及可数集 $\mathcal X_1(L^{(i)})$）。把那里确定的
多值截面记作 $\mathfrak s_{\beta^{(i)},\vec P^{(i)}}$。

\begin{proposition}{3.7.36}\label{prop:3-7-36}
在上述情形中，模空间
\[
 \mathcal M_{k_1,k_0}
 (L^{(1)},L^{(0)};[\ell_p,w_1],[\ell_q,w_2];
 \vec P^{(1)},\vec P^{(0)})
\]
上存在多值截面系
\[
 \mathfrak s_{[\ell_p,w_1],[\ell_q,w_2],
 \vec P^{(1)},\vec P^{(0)}},
\]
并具有下列性质：
\begin{equation}
 \text{上述模空间的虚维数由式 \eqref{eq:3-7-38} 给出。}
\tag{3.7.37.1}\label{eq:3-7-37-1}
\end{equation}
\begin{equation}
\begin{aligned}
 &\mu([\ell_q,w_2])-\mu([\ell_p,w_1])
 -\sum_{i=1}^{k_0}\deg P_i^{(0)}
 -\sum_{i=1}^{k_1}\deg P_i^{(1)}
 +k_0+k_1-1.
\end{aligned}
\tag{3.7.38}\label{eq:3-7-38}
\end{equation}
这里 $\deg P_i^{(0)}$ 与 $\deg P_i^{(1)}$ 是移位前把
$P_i^{(0)},P_i^{(1)}$ 看作上链时的次数。
\begin{equation}
 \mathfrak s_{[\ell_p,w_1],[\ell_q,w_2],
 \vec P^{(1)},\vec P^{(0)}}\text{ 横截于 }0.
\tag{3.7.37.2}\label{eq:3-7-37-2}
\end{equation}
\begin{equation}
\begin{gathered}
 \mathfrak s_{[\ell_p,w_1],[\ell_q,w_2],
 \vec P^{(1)},\vec P^{(0)}}\text{ 在边界处与其他同类多值截面}\\
 \text{以及 }
 \mathfrak s_{\beta^{(0)},\vec P^{(0)}},
 \mathfrak s_{\beta^{(1)},\vec P^{(1)}}\text{ 相容}.
\end{gathered}
\tag{3.7.37.3}\label{eq:3-7-37-3}
\end{equation}
\end{proposition}

\begin{proof}
证明与滤过 $A_\infty$ 代数的情形（命题 3.5.2）相同，细节见第 7.2 节。
维数可仿照第 2.3 节及命题 3.5.2 的方法计算。
\end{proof}

\begin{definition}{3.7.39}\label{def:3-7-39}
记
\[
\begin{aligned}
 &\mathcal M_{k_1,k_0}
 (L^{(1)},L^{(0)};[\ell_p,w_1],[\ell_q,w_2];
 \vec P^{(1)},\vec P^{(0)})^{\mathfrak s}\\
 &\quad:=
 \left(
 \mathfrak s_{[\ell_p,w_1],[\ell_q,w_2],
 \vec P^{(1)},\vec P^{(0)}}
 \right)^{-1}(0).
\end{aligned}
\]
\end{definition}

上述模空间的定向见第 8 章定义 8.7.1。它是一个带三角剖分的空间，并且每个
最高维单形都带有有理权重（见 \textbf{[FuOn99II]} 第 6 节或定义 A1.28）。
若 \eqref{eq:3-7-38} 的维数为 $0$，它由有限多个点组成，每个点都带一个
$\mathbb Q$ 中的权重；若维数为 $1$，它是一张定向图，每条边都带一个
$\mathbb Q$ 中的权重。
\liftIndexRangeEnd{I126}

\begin{remark}{3.7.40}\label{rem:3-7-40}
本小节没有在扰动后的模空间
\[
 \mathcal M_{k_1,k_0}
 (L^{(1)},L^{(0)};[\ell_p,w_1],[\ell_q,w_2];
 \vec P^{(1)},\vec P^{(0)})^{\mathfrak s}
\]
上指定任何评价映射，所以这里不把它视为奇异链。不过，更自然的做法是把它看作
零维空间 $\{q\}\subset L^{(1)}\cap L^{(0)}$ 上的奇异链。下一小节处理
$L^{(1)}\cap L^{(0)}$ 是有限维流形之并的情形时，这一点会更加清楚。
\end{remark}

现在利用这些模空间定义运算族 $\mathfrak n_{k_1,k_0}$，它们将在
$C(L^{(1)},L^{(0)};\Lambda_{0,\mathrm{nov}})$ 上给出滤过
$A_\infty$ 双模结构。

\begin{definition}{3.7.41}\label{def:3-7-41}
考虑 \eqref{eq:3-7-38} 的维数为 $0$ 的情形。设
\[
\begin{gathered}
 T^{\lambda_i^{(0)}}e^{\mu_i^{(0)}}P_i^{(0)}
 \in C(L^{(0)};\Lambda_{0,\mathrm{nov}}),\\
 T^{\lambda_i^{(1)}}e^{\mu_i^{(1)}}P_i^{(1)}
 \in C(L^{(1)};\Lambda_{0,\mathrm{nov}}),
\end{gathered}
\]
并设
\[
 T^\lambda e^\mu[\ell_p,w_1]
 \in C(L^{(1)},L^{(0)};\Lambda_{0,\mathrm{nov}}).
\]
定义
\[
\begin{aligned}
&\mathfrak n_{k_1,k_0}\bigl(
 T^{\lambda_1^{(1)}}e^{\mu_1^{(1)}}P_1^{(1)}
 \otimes\cdots\otimes
 T^{\lambda_{k_1}^{(1)}}e^{\mu_{k_1}^{(1)}}P_{k_1}^{(1)}
 \otimes T^\lambda e^\mu[\ell_p,w_1]\\
&\hspace{8em}\otimes
 T^{\lambda_1^{(0)}}e^{\mu_1^{(0)}}P_1^{(0)}
 \otimes\cdots\otimes
 T^{\lambda_{k_0}^{(0)}}e^{\mu_{k_0}^{(0)}}P_{k_0}^{(0)}
 \bigr)\\
&\quad:=
 \sum_{[\ell_q,w_2]\in\operatorname{Cr}(L^{(0)},L^{(1)})}
 \#\mathcal M_{k_1,k_0}
 (L^{(1)},L^{(0)};[\ell_p,w_1],[\ell_q,w_2];
 \vec P^{(1)},\vec P^{(0)})^{\mathfrak s}\\
&\hspace{13em}\cdot
 T^{\lambda'}e^{\mu'}[\ell_q,w_2],
\end{aligned}
\]
其中
\[
 \lambda'
 =\sum_{i=1}^{k_0}\lambda_i^{(0)}
  +\lambda+\sum_{i=1}^{k_1}\lambda_i^{(1)},
 \qquad
 \mu'
 =\sum_{i=1}^{k_0}\mu_i^{(0)}
  +\mu+\sum_{i=1}^{k_1}\mu_i^{(1)}.
\]
这里 $\#$ 表示各点所带权重之和。
\end{definition}

若伪全纯映射 $u\colon\Sigma\to M$ 满足
\eqref{eq:3-7-22-1}--\eqref{eq:3-7-22-6}，则由
\eqref{eq:3-7-22-5}，
\[
 \int_{w_2}\omega-\int_{w_1}\omega
 =\int_\Sigma u^*\omega\geq0.
\]
因此，$T^{\lambda'}e^{\mu'}[\ell_q,w_2]$ 的能量满足
\begin{equation}
\begin{aligned}
 \lambda'+\int_{w_2}\omega
 \geq{}&
 \left(\sum_{i=1}^{k_0}\lambda_i^{(0)}\right)
 +\left(\sum_{i=1}^{k_1}\lambda_i^{(1)}\right)
 +\left(\lambda+\int_{w_1}\omega\right).
\end{aligned}
\tag{3.7.42}\label{eq:3-7-42}
\end{equation}
右端每一项都非负，故所得元素确实属于
$C(L^{(1)},L^{(0)};\Lambda_{0,\mathrm{nov}})$。于是
\[
\begin{aligned}
\mathfrak n_{k_1,k_0}\colon{}&
B_{k_1}\bigl(C(L^{(1)};\Lambda_{0,\mathrm{nov}})[1]\bigr)
\otimes_{\Lambda_{0,\mathrm{nov}}}
C(L^{(1)},L^{(0)};\Lambda_{0,\mathrm{nov}})[1]\\
&\otimes_{\Lambda_{0,\mathrm{nov}}}
B_{k_0}\bigl(C(L^{(0)};\Lambda_{0,\mathrm{nov}})[1]\bigr)\\
&\longrightarrow
C(L^{(1)},L^{(0)};\Lambda_{0,\mathrm{nov}})[1].
\end{aligned}
\]
定义 3.7.41 中只对 \eqref{eq:3-7-38} 等于 $0$ 的
$[\ell_q,w_2]$ 求和，因此 $\mathfrak n_{k_1,k_0}$ 的次数为 $+1$。
此外，由 (3.2.12.1)，不等式 \eqref{eq:3-7-42} 蕴含这些映射满足条件
\eqref{eq:3-7-1}。

下面验证有隙条件。以
\[
 G(L^{(1)},L^{(0)})_0
 \subset\mathbb R_{\geq0}\times2\mathbb Z
\]
表示由下列集合生成的子幺半群：
\[
 \left\{
  (\omega[B],\mu(p,q;B))
  \ \middle|\
  \begin{gathered}
   p,q\in L^{(0)}\cap L^{(1)},\quad B\in\pi_2(p,q),\\
   \mathcal M(p,q;B)\ne\varnothing
  \end{gathered}
 \right\}.
\]
记号见第 2.3.1 小节。定义\liftIndex{I128}
\begin{equation}
\begin{aligned}
 G(L^{(1)},L^{(0)})
 :=\{\,\beta^{(1)}+\alpha+\beta^{(0)}
 \mid{}&
 \beta^{(i)}\in G(L^{(i)}),\\
 &\alpha\in G(L^{(1)},L^{(0)})_0\,\}.
\end{aligned}
\tag{3.7.43}\label{eq:3-7-43}
\end{equation}
它是包含 $G(L^{(1)})$ 与 $G(L^{(0)})$ 的子幺半群；由格罗莫夫紧性定理\liftIndex{I232}，
$G(L^{(1)},L^{(0)})$ 满足条件 3.1.6。按定义，我们的运算是
$G(L^{(1)},L^{(0)})$-有隙的。

现在令 $\widehat d$ 为由 $\mathfrak n_{k_1,k_0}$ 依照
\eqref{eq:3-7-2} 得到的映射。要证明定理 3.7.21，只需证明下面的结论。
\translatornote{原书此处把运算族误写成 $\mathfrak n_{k_1,k_2}$；
本节双模运算的第二个指标始终是 $k_0$，译文据此改正。}

\begin{lemma}{3.7.44}\label{lem:3-7-44}
由 \eqref{eq:3-7-2} 定义的同态 $\widehat d$ 满足
\[
 \widehat d\circ\widehat d=0.
\]
\end{lemma}

\begin{proof}[引理 3.7.44 的证明]
考虑维数为 $1$ 的模空间
\[
 \mathcal M_{k_1,k_0}
 (L^{(1)},L^{(0)};[\ell_p,w_1],[\ell_q,w_2];
 \vec P^{(1)},\vec P^{(0)})^{\mathfrak s}.
\]
它的边界由图 3.7.5--3.7.7 中三种情形描述。

\begin{figure}[htbp]
  \centering
  \setlength{\unitlength}{0.88mm}
  \begin{picture}(112,58)
    \qbezier(7,35)(55,65)(105,35)
    \qbezier(7,35)(55,5)(105,35)
    \put(3,31){\line(1,1){7}}
    \put(3,39){\line(1,-1){7}}
    \put(102,31){\line(1,1){7}}
    \put(102,39){\line(1,-1){7}}
    \put(32,24){\line(1,2){11}}
    \put(53,20){\line(1,2){14}}
    \put(75,23){\line(1,2){11}}
    \put(59,8){\circle{18}}
    \put(59,17){\circle*{1.5}}
    \put(53,3){\line(1,1){3}}
    \put(53,6){\line(1,-1){3}}
    \put(64,3){\line(1,1){3}}
    \put(64,6){\line(1,-1){3}}
    \put(1,25){$p$}
    \put(106,25){$q$}
    \put(20,48){$P_{k_1}^{(1)}$}
    \put(87,48){$P_1^{(1)}$}
    \put(17,15){$P_1^{(0)}$}
    \put(88,16){$P_{k_0}^{(0)}$}
    \put(43,-2){$P_*^{(0)}$}
    \put(66,-2){$P_*^{(0)}$}
    \put(106,43){$L^{(0)}$}
    \put(106,18){$L^{(1)}$}
  \end{picture}
  \par\smallskip
  \renewcommand{\thefigure}{3.7.5}
  \refstepcounter{figure}
  \textbf{\figurename~\thefigure}\label{fig:3-7-5}
\end{figure}

\begin{figure}[htbp]
  \centering
  \setlength{\unitlength}{0.88mm}
  \begin{picture}(112,58)
    \qbezier(7,25)(55,55)(105,25)
    \qbezier(7,25)(55,-5)(105,25)
    \put(3,21){\line(1,1){7}}
    \put(3,29){\line(1,-1){7}}
    \put(102,21){\line(1,1){7}}
    \put(102,29){\line(1,-1){7}}
    \put(32,14){\line(1,2){11}}
    \put(53,10){\line(1,2){14}}
    \put(75,13){\line(1,2){11}}
    \put(57,49){\circle{18}}
    \put(57,40){\circle*{1.5}}
    \put(51,53){\line(1,1){3}}
    \put(51,56){\line(1,-1){3}}
    \put(62,53){\line(1,1){3}}
    \put(62,56){\line(1,-1){3}}
    \put(1,15){$p$}
    \put(106,15){$q$}
    \put(18,39){$P_{k_1}^{(1)}$}
    \put(87,39){$P_1^{(1)}$}
    \put(17,5){$P_1^{(0)}$}
    \put(88,6){$P_{k_0}^{(0)}$}
    \put(39,55){$P_*^{(1)}$}
    \put(67,55){$P_*^{(1)}$}
    \put(106,33){$L^{(0)}$}
    \put(106,8){$L^{(1)}$}
  \end{picture}
  \par\smallskip
  \renewcommand{\thefigure}{3.7.6}
  \refstepcounter{figure}
  \textbf{\figurename~\thefigure}\label{fig:3-7-6}
\end{figure}

\begin{figure}[htbp]
  \centering
  \setlength{\unitlength}{0.88mm}
  \begin{picture}(128,47)
    \qbezier(5,24)(34,46)(63,24)
    \qbezier(5,24)(34,2)(63,24)
    \qbezier(63,24)(94,47)(124,24)
    \qbezier(63,24)(94,1)(124,24)
    \put(1,20){\line(1,1){7}}
    \put(1,28){\line(1,-1){7}}
    \put(121,20){\line(1,1){7}}
    \put(121,28){\line(1,-1){7}}
    \put(62.2,23.2){\circle*{1.8}}
    \put(20,16){\line(1,2){9}}
    \put(38,14){\line(1,2){9}}
    \put(75,17){\line(1,2){8}}
    \put(99,14){\line(1,2){9}}
    \put(0,12){$p$}
    \put(124,12){$q$}
    \put(0,29.5){$L^{(0)}$}
    \put(0,6.5){$L^{(1)}$}
    \put(117,29.5){$L^{(1)}$}
    \put(106,10.4){$L^{(0)}$}
    \put(16,39){$P_{k_1}^{(1)}$}
    \put(107,39){$P_1^{(1)}$}
    \put(15,4){$P_1^{(0)}$}
    \put(105,5){$P_{k_0}^{(0)}$}
    \put(48,34){$P_*^{(1)}$}
    \put(67,34){$P_*^{(1)}$}
    \put(48,8){$P_*^{(0)}$}
    \put(67,8){$P_*^{(0)}$}
  \end{picture}
  \par\smallskip
  \renewcommand{\thefigure}{3.7.7}
  \refstepcounter{figure}
  \textbf{\figurename~\thefigure}\label{fig:3-7-7}
\end{figure}

图 3.7.5 对应的各项为
\[
 \mathfrak n\left(
 \widehat d^{(1)}
 (P_1^{(1)}\otimes\cdots\otimes P_{k_1}^{(1)})
 \otimes[\ell_p,w_1]\otimes
 P_1^{(0)}\otimes\cdots\otimes P_{k_0}^{(0)}
 \right),
\]
其中 $\widehat d^{(1)}$ 定义于
$\bigl(C(L^{(1)};\Lambda_{0,\mathrm{nov}}),\mathfrak m^{(1)}\bigr)$。
图 3.7.6 对应的各项为
\[
 \mathfrak n\left(
 P_1^{(1)}\otimes\cdots\otimes P_{k_1}^{(1)}
 \otimes[\ell_p,w_1]\otimes
 \widehat d^{(0)}
 (P_1^{(0)}\otimes\cdots\otimes P_{k_0}^{(0)})
 \right),
\]
其中 $\widehat d^{(0)}$ 定义于
$\bigl(C(L^{(0)};\Lambda_{0,\mathrm{nov}}),\mathfrak m^{(0)}\bigr)$。
图 3.7.7 对应的各项为
\[
 \mathfrak n\left(
 P_1^{(1)}\otimes\cdots\otimes
 \mathfrak n(\cdots\otimes P_{k_1}^{(1)}
 \otimes[\ell_p,w_1]\otimes P_1^{(0)}\otimes\cdots)
 \otimes\cdots\otimes P_{k_0}^{(0)}
 \right).
\]
合并第 8.5、8.7 节的论证，可以用类似方式核对符号；上述三类项彼此抵消，
从而证明引理 3.7.44。
\end{proof}

至此，我们在 $C(L^{(1)},L^{(0)};\Lambda_{0,\mathrm{nov}})$ 上构造了左
$\bigl(C(L^{(1)};\Lambda_{0,\mathrm{nov}}),\mathfrak m^{(1)}\bigr)$、
右 $\bigl(C(L^{(0)};\Lambda_{0,\mathrm{nov}}),\mathfrak m^{(0)}\bigr)$
的滤过 $A_\infty$ 双模结构，定理 3.7.21 得证。
\end{proof}

现在假设
$\bigl(C(L^{(0)};\Lambda_{0,\mathrm{nov}}),\mathfrak m^{(0)}\bigr)$
与
$\bigl(C(L^{(1)};\Lambda_{0,\mathrm{nov}}),\mathfrak m^{(1)}\bigr)$
都无障碍。由上一小节的定义--引理 3.7.13 和引理 3.7.14，对任意定界余链
\[
 b_i\in\widehat{\mathcal M}(L^{(i)};J_i;\mathfrak s_i)
 \qquad(i=0,1)
\]
可以定义上边界算子
\[
 \delta_{b_1,b_0}\colon
 C(L^{(1)},L^{(0)};\Lambda_{0,\mathrm{nov}})
 \longrightarrow
 C(L^{(1)},L^{(0)};\Lambda_{0,\mathrm{nov}}).
\]
这给出了标准弗洛尔上边界算子的形变。因此，下面的定理是引理 3.7.14 和
定理 3.7.21 的直接推论。

\begin{theorem}{3.7.45}\label{thm:3-7-45}
设 $(L^{(0)},L^{(1)})$ 是一对彼此横截的相对 spin 拉格朗日子流形，并设
\[
 b_0\in\widehat{\mathcal M}(L^{(0)};J_0;\mathfrak s_0),
 \qquad
 b_1\in\widehat{\mathcal M}(L^{(1)};J_1;\mathfrak s_1).
\]
则存在上同调
\[
 HF\bigl((L^{(1)},b_1),(L^{(0)},b_0);
 \Lambda_{0,\mathrm{nov}}\bigr)
 =
 \frac{\operatorname{Ker}\delta_{b_1,b_0}}
      {\operatorname{Im}\delta_{b_1,b_0}},
\]
称为由 $b_0,b_1$ 形变的 $(L^{(0)},L^{(1)})$ 的弗洛尔上同调
\liftIndex{I125}。
\end{theorem}

当然，如果上述两个滤过 $A_\infty$ 代数关于同伦幺元弱无障碍，则可以应用
命题 3.7.17：只要
\[
 b_i\in\widehat{\mathcal M}_{\mathrm{weak}}
 (L^{(i)};J_i;\mathfrak s_i),
 \qquad
 \mathfrak{PO}_0(b_0)=\mathfrak{PO}_1(b_1),
\]
便可定义由 $b_i$ 形变的弗洛尔上同调。

\begin{remark}{3.7.46}\label{rem:3-7-46}
仿照定义 3.7.20，定义
\[
 C(L^{(1)},L^{(0)};\ell_0;\Lambda_{\mathrm{nov}})
 :=
 \bigl(
 C(L^{(1)},L^{(0)};\ell_0;\mathbb Q)
 \otimes\Lambda_{\mathrm{nov}}
 \bigr)/\sim
\]
并置
\[
 C(L^{(1)},L^{(0)};\ell_0;\Lambda_{0,\mathrm{nov}})
 :=
 F^0C(L^{(1)},L^{(0)};\ell_0;\Lambda_{\mathrm{nov}}).
\]
由构造容易看出，分解
\[
 C(L^{(1)},L^{(0)};\Lambda_{0,\mathrm{nov}})
 =
 \widehat{\bigoplus}_{
 [\ell_0]\in\pi_0(\Omega(L^{(0)},L^{(1)}))}
 C(L^{(1)},L^{(0)};\ell_0;\Lambda_{0,\mathrm{nov}})
\]
是滤过 $A_\infty$ 双模的完备直和。因而弗洛尔上同调也分解为
\begin{equation}
\begin{aligned}
&HF\bigl((L^{(1)},b_1),(L^{(0)},b_0);
 \Lambda_{0,\mathrm{nov}}\bigr)\\
&\quad\cong
 \widehat{\bigoplus}_{
 [\ell_0]\in\pi_0(\Omega(L^{(0)},L^{(1)}))}
 HF\bigl((L^{(1)},b_1),(L^{(0)},b_0);
 \ell_0;\Lambda_{0,\mathrm{nov}}\bigr).
\end{aligned}
\tag{3.7.47}\label{eq:3-7-47}
\end{equation}
下一小节讨论的博特--莫尔斯情形也适用同一说明。
\end{remark}
